\NormalFont\ShortTitle{Evrei}

{\MT Scrisoarea către Evrei


\par }\ChapOne{1}{\PP \VerseOne{1}Dumnezeu, care în trecut a vorbit părinților, prin prooroci, în multe vremuri și în felurite moduri,

\VS{2}la sfârșitul acestor zile, ne-a vorbit prin Fiul Său, pe care L-a pus moștenitor al tuturor lucrurilor și prin care a făcut și lumile.

\VS{3}Fiul său este strălucirea gloriei sale, chipul însuși al substanței sale și care susține toate lucrurile prin cuvântul puterii sale, care, după ce ne-a purificat prin el însuși de păcatele noastre, a șezut la dreapta Majestății de sus,

\VS{4}devenind cu atât mai bun decât îngerii, cu cât numele mai presus de îngerii pe care l-a moștenit este mai presus decât al lor.

\VS{5}Căci cui dintre îngeri i-a spus el vreodată

\par }{\Q “Tu ești Fiul meu.

\par }{\QB Azi am devenit tatăl tău?”

\par }{\MM și din nou,

\par }{\Q “Îi voi fi un Tată,

\par }{\QB Și el îmi va fi pentru mine un fiu?”


\par }{\PP \VS{6}Când va aduce din nou pe lume pe întâiul născut, zice: “Toți îngerii lui Dumnezeu să i se închine”.

\VS{7}Despre îngeri spune,

\par }{\Q “El își face îngerii vânturi,

\par }{\QB și robii lui o flacără de foc.”


\par }{\PP \VS{8}Dar despre Fiul spune,

\par }{\Q “Tronul Tău, Dumnezeule, este în vecii vecilor.

\par }{\QB Sceptrul dreptății este sceptrul Împărăției Tale.


\par }{\Q \VS{9}Tu ai iubit dreptatea și ai urât nelegiuirea;

\par }{\QB De aceea Dumnezeu, Dumnezeul tău, te-a uns cu untdelemnul bucuriei mai presus de semenii tăi.”


\par }{\PP \VS{10}Și,

\par }{\Q “Tu, Doamne, la început ai pus temelia pământului.

\par }{\QB Cerurile sunt opera mâinilor Tale.


\par }{\Q \VS{11}Ei vor pieri, dar tu vei rămâne.

\par }{\QB Toate vor îmbătrâni ca o haină.


\par }{\Q \VS{12}Le vei înfășura ca pe o mantie,

\par }{\QB și vor fi schimbate;

\par }{\QB dar tu ești la fel.

\par }{\QB Anii tăi nu vor da greș.”


\par }{\PP \VS{13}Dar care dintre îngeri a spus vreodată,

\par }{\Q “Stai la dreapta mea,

\par }{\QB până ce voi face din vrăjmașii tăi scăunelul picioarelor tale?”


\par }{\PP \VS{14}Nu sunt oare toți duhuri slujitoare, trimise să slujească în folosul celor ce vor moșteni mântuirea?



\par }\Chap{2}{\PP \VerseOne{1}De aceea trebuie să fim mai atenți la cele auzite, ca nu cumva să ne îndepărtăm.

\VS{2}Căci, dacă cuvântul rostit prin îngeri s-a dovedit statornic, iar orice fărădelege și neascultare a primit o pedeapsă dreaptă,

\VS{3}cum vom scăpa dacă vom neglija o mântuire atât de mare — care, la început, după ce a fost rostită prin Domnul, ne-a fost confirmată de cei care au auzit-o,

\VS{4}Dumnezeu mărturisind și el împreună cu ei, atât prin semne și minuni, cât și prin diferite lucrări de putere și prin daruri ale Duhului Sfânt, după voia lui?


\par }{\PP \VS{5}Căci El nu a supus lumii viitoare, despre care vorbim, îngerilor.

\VS{6}Ci cineva a mărturisit undeva, zicând:,

\par }{\Q “Ce este omul, de te gândești la el?

\par }{\QB Sau pe fiul omului, ca să vă pese de el?


\par }{\Q \VS{7}L-ai făcut puțin mai prejos decât îngerii.

\par }{\QB L-ai încoronat cu glorie și onoare.


\par }{\QB \VS{8}Tu ai supus toate lucrurile sub picioarele Lui.”

\par }{\PP Pentru că, supunându-i toate lucrurile, nu a lăsat nimic din ceea ce nu-i este supus lui. Dar acum nu vedem încă toate lucrurile supuse lui.

\VS{9}Ci îl vedem pe cel care a fost făcut cu puțin mai prejos decât îngerii, pe Isus, din cauza suferinței morții, încoronat cu glorie și onoare, pentru ca, prin harul lui Dumnezeu, să guste din moarte pentru toți.


\par }{\PP \VS{10}Fiindcă Lui, pentru care sunt toate lucrurile și prin care sunt toate lucrurile, i s-a cuvenit, aducând mulți copii la slavă, să desăvârșească prin suferințe pe autorul mântuirii lor.

\VS{11}Căci atât cel ce sfințește, cât și cei ce sunt sfințiți sunt toți de la unul singur, motiv pentru care nu se rușinează să-i numească frați,

\VS{12}spunând:

\par }{\Q “Voi vesti numele tău fraților mei.

\par }{\QB În mijlocul adunării voi cânta lauda ta.”


\par }{\PP \VS{13}Și iarăși: “În El îmi voi pune încrederea mea.” Din nou: “Iată-mă aici cu copiii pe care mi i-a dat Dumnezeu.”

\VS{14}De vreme ce copiii s-au împărtășit din carne și sânge, și el însuși s-a împărtășit în același fel, pentru ca, prin moarte, să nimicească pe cel care avea puterea morții, adică pe diavolul,

\VS{15}și să elibereze pe toți cei care, de frica morții, au fost toată viața lor supuși robiei.

\VS{16}Pentru că, foarte sigur, nu dă ajutor îngerilor, ci dă ajutor urmașilor lui Avraam.

\VS{17}De aceea a fost obligat în toate lucrurile să se facă asemenea fraților săi, ca să devină un mare preot milostiv și credincios în ceea ce privește lucrurile care țin de Dumnezeu, pentru a face ispășire pentru păcatele poporului.

\VS{18}Căci, întrucât el însuși a suferit fiind ispitit, poate ajuta pe cei ce sunt ispitiți.



\par }\Chap{3}{\PP \VerseOne{1}Așadar, frați sfinți, părtași ai chemării cerești, luați aminte la Apostolul și Marele Preot al mărturisirii noastre: Isus,

\VS{2}care a fost credincios celui care l-a rânduit, cum a fost și Moise în toată casa lui.

\VS{3}Căci el a fost socotit vrednic de mai multă glorie decât Moise, pentru că cel care a zidit casa are mai multă cinste decât casa.

\VS{4}Căci orice casă este construită de cineva, dar cel care a construit totul este Dumnezeu.

\VS{5}Într-adevăr, Moise a fost credincios în toată casa sa ca un slujitor, ca mărturie a lucrurilor care urmau să fie spuse mai târziu,

\VS{6}dar Cristos este credincios ca un Fiu peste casa sa. Noi suntem casa lui, dacă ne păstrăm până la sfârșit încrederea și slava speranței noastre fermă.

\VS{7}De aceea, așa cum spune Duhul Sfânt,

\par }{\Q “Astăzi, dacă veți auzi vocea lui,


\par }{\QB \VS{8}Nu vă împietriți inimile ca la răzvrătire,

\par }{\QB în ziua încercării în pustie,


\par }{\QB \VS{9}unde m-au încercat și m-au pus la încercare părinții voștri,

\par }{\QB și a văzut faptele mele timp de patruzeci de ani.


\par }{\Q \VS{10}De aceea am fost nemulțumit de neamul acela,

\par }{\QB ș i a spus: “Ei greș esc întotdeauna în inima lor,

\par }{\QB dar ei nu-mi cunoșteau căile.


\par }{\Q \VS{11}Așa cum am jurat în mânia Mea,

\par }{\QB “Nu vor intra în odihna Mea.”


\par }{\PP \VS{12}Luați seama, fraților, ca nu cumva să fie în vreunul din voi o inimă rea, necredincioasă, care să se depărteze de Dumnezeul cel viu;

\VS{13}ci îndemnați-vă unii pe alții în fiecare zi, atâta timp cât se numește “astăzi”, ca nu cumva vreunul din voi să se împietrească prin înșelăciunea păcatului.

\VS{14}Căci am devenit părtași ai lui Hristos, dacă ne păstrăm până la sfârșit începutul încrederii noastre,

\VS{15}în timp ce se spune

\par }{\Q “Astăzi, dacă veți auzi vocea lui,

\par }{\QB nu vă împietriți inimile, ca în timpul răzvrătirii.”


\par }{\PP \VS{16}Căci cine s-a răzvrătit când a auzit? Nu au fost oare toți cei care au ieșit din Egipt conduși de Moise?

\VS{17}Cu cine a fost nemulțumit patruzeci de ani? Nu a fost oare cu cei care au păcătuit, ale căror trupuri au căzut în pustiu?

\VS{18}Cui a jurat că nu vor intra în odihna lui, ci celor neascultători?

\VS{19}Vedem că ei nu au putut intra din cauza necredinței.



\par }\Chap{4}{\PP \VerseOne{1}Să ne temem, așadar, ca nu cumva vreunul dintre voi să pară că a lipsit de promisiunea de a intra în odihna Lui.

\VS{2}Căci, într-adevăr, ni s-a vestit o veste bună, așa cum au făcut și ei, dar cuvântul pe care l-au auzit nu le-a fost de folos, pentru că nu a fost amestecat cu credință de către cei care l-au auzit.

\VS{3}Căci noi, cei care am crezut, intrăm în acea odihnă, așa cum a spus el: “Așa cum am jurat în mânia Mea, nu vor intra în odihna Mea”, deși lucrările au fost terminate de la întemeierea lumii.

\VS{4}Căci el a spus undeva despre ziua a șaptea: “În ziua a șaptea, Dumnezeu s-a odihnit în ziua a șaptea de toate lucrările Sale”;

\VS{5}iar în acest loc din nou: “Ei nu vor intra în odihna Mea”.


\par }{\PP \VS{6}Văzând, așadar, că rămâne ca unii să intre în ea, iar cei cărora le-a fost vestită mai înainte vestea cea bună nu au reușit să intre din cauza neascultării,

\VS{7}el definește din nou o zi anume, “astăzi”, spunând prin David, la o distanță atât de mare de timp după aceea (așa cum s-a spus),

\par }{\Q “Astăzi, dacă veți auzi vocea lui,

\par }{\QB nu vă împietriți inimile.”


\par }{\PP \VS{8}Căci dacă Iosua le-ar fi dat odihnă, n-ar fi vorbit după aceea de o altă zi.

\VS{9}Rămâne așadar o odihnă de Sabat pentru poporul lui Dumnezeu.

\VS{10}Căci cel care a intrat în odihna lui s-a odihnit și el însuși de lucrările sale, cum s-a odihnit și Dumnezeu de ale lui.

\VS{11}Să ne sârguim deci să intrăm în această odihnă, ca nu cumva cineva să cadă după același exemplu de neascultare.

\VS{12}Căci Cuvântul lui Dumnezeu este viu și lucrător și mai ascuțit decât orice sabie cu două tăișuri, străpungând până la despărțirea sufletului și a duhului, a încheieturilor și a măduvei, și poate să discearnă gândurile și intențiile inimii.

\VS{13}Nicio făptură nu este ascunsă de ochii Lui, ci toate lucrurile sunt goale și expuse înaintea ochilor Celui căruia trebuie să dăm socoteală.


\par }{\PP \VS{14}Deci, având un mare preot care a trecut prin ceruri, pe Isus, Fiul lui Dumnezeu, să ne ținem cu tărie mărturisirea noastră.

\VS{15}Căci noi nu avem un mare preot care nu poate fi atins de sentimentul slăbiciunilor noastre, ci unul care a fost în toate privințele ispitit ca și noi, dar fără păcat.

\VS{16}Să ne apropiem, așadar, cu îndrăzneală de tronul harului, ca să primim îndurare și să găsim har pentru ajutor la vreme de nevoie.



\par }\Chap{5}{\PP \VerseOne{1}Fiindcă orice mare preot, luat dintre oameni, este rânduit pentru oameni în ceea ce privește lucrurile care țin de Dumnezeu, ca să aducă daruri și jertfe pentru păcate.

\VS{2}Marele preot poate să se poarte cu blândețe cu cei care sunt neștiutori și se rătăcesc, pentru că și el însuși este înconjurat de slăbiciune.

\VS{3}Din această cauză, el trebuie să ofere jertfe pentru păcate atât pentru popor, cât și pentru el însuși.

\VS{4}Nimeni nu-și asumă singur această onoare, ci este chemat de Dumnezeu, așa cum a fost Aaron.

\VS{5}Tot astfel, nici Hristos nu s-a proslăvit pe sine însuși pentru a fi făcut mare preot, ci el a fost cel care i-a spus

\par }{\Q “Tu ești Fiul meu.

\par }{\QB Astăzi am devenit tatăl tău.”


\par }{\PP \VS{6}După cum spune și în alt loc,

\par }{\Q “Ești preot pentru totdeauna,

\par }{\QB după ordinul lui Melchisedec.”


\par }{\PP \VS{7}El, în zilele trupului său, după ce a făcut rugăciuni și cereri cu strigăte și lacrimi puternice către Cel ce putea să-l scape de la moarte, și după ce a fost ascultat pentru frica lui evlavioasă,

\VS{8}deși era un Fiu, a învățat ascultarea prin lucrurile pe care le-a suferit.

\VS{9}Fiind desăvârșit, el a devenit pentru toți cei care îl ascultă autorul mântuirii veșnice,

\VS{10}fiind numit de Dumnezeu mare preot după ordinul lui Melchisedec.


\par }{\PP \VS{11}Despre El avem multe cuvinte de spus și greu de înțeles, pentru că v-ați pierdut auzul.

\VS{12}Căci, deși până acum ar trebui să fiți învățători, aveți din nou nevoie ca cineva să vă învețe rudimentele primelor principii ale revelațiilor lui Dumnezeu. Ați ajuns să aveți nevoie de lapte, și nu de hrană solidă.

\VS{13}Căci oricine trăiește cu lapte nu are experiență în cuvântul dreptății, pentru că este un copil.

\VS{14}Dar hrana solidă este pentru cei maturi, care, prin folosire, și-au exersat simțurile pentru a discerne binele și răul.



\par }\Chap{6}{\PP \VerseOne{1}Așadar, lăsând învățătura primelor principii ale lui Hristos, să mergem mai departe spre desăvârșire — să nu punem din nou temelia pocăinței de faptele moarte, a credinței în Dumnezeu,

\VS{2}a învățăturii despre botezuri, despre punerea mâinilor, despre învierea morților și despre judecata veșnică.

\VS{3}Asta vom face, dacă Dumnezeu ne va îngădui.

\VS{4}Căci în ceea ce privește pe cei care au fost odată luminați și au gustat din darul ceresc și au fost făcuți părtași la Duhul Sfânt,

\VS{5}și au gustat cuvântul bun al lui Dumnezeu și puterile veacului viitor,

\VS{6}și apoi au căzut, este imposibil să-i reînnoim din nou la pocăință, întrucât ei îl răstignesc din nou pe Fiul lui Dumnezeu pentru ei înșiși și îl fac de rușine.

\VS{7}Căci pământul care a băut ploaia care vine des peste el și produce o recoltă potrivită pentru cei pentru care este și lucrat, primește binecuvântare de la Dumnezeu;

\VS{8}dar dacă poartă spini și ciulini, este respins și aproape de a fi blestemat, al cărui sfârșit este să fie ars.


\par }{\PP \VS{9}Dar, preaiubiților, noi suntem încredințați de lucruri mai bune pentru voi și de lucruri care însoțesc mântuirea, chiar dacă vorbim astfel.

\VS{10}Căci Dumnezeu nu este nedrept, ca să uite lucrarea voastră și osteneala dragostei pe care ați arătat-o față de numele Lui, prin faptul că ați slujit sfinților și încă le mai slujiți.

\VS{11}Dorim ca fiecare dintre voi să dea dovadă de aceeași sârguință în vederea împlinirii speranței până la sfârșit,

\VS{12}ca să nu fiți leneși, ci imitatori ai celor care, prin credință și perseverență, au moștenit promisiunile.


\par }{\PP \VS{13}Căci, când Dumnezeu a făcut făgăduința lui Avraam, n-a putut să jure pe nimeni mai mare, ci a jurat pe sine însuși,

\VS{14}zicând: “Te voi binecuvânta și te voi binecuvânta și te voi înmulți”.

\VS{15}Astfel, după ce a răbdat, a obținut promisiunea.

\VS{16}Căci, într-adevăr, oamenii jură pe unul mai mare și, în orice dispută a lor, jurământul este definitiv pentru confirmare.

\VS{17}În felul acesta, Dumnezeu, hotărât să arate mai mult moștenitorilor promisiunii imuabilitatea sfatului său, a intervenit cu un jurământ,

\VS{18}pentru ca, prin două lucruri imuabile, în care este imposibil ca Dumnezeu să mintă, să avem o puternică încurajare, noi, care am fugit să ne refugiem pentru a ne agăța de speranța pusă înaintea noastră.

\VS{19}Această nădejde o avem ca o ancoră a sufletului, o nădejde deopotrivă sigură și neclintită și care intră în ceea ce este înăuntrul perdelei,

\VS{20}unde, ca un precursor, Isus a intrat pentru noi, devenind Mare Preot pentru totdeauna, după rânduiala lui Melchisedec.



\par }\Chap{7}{\PP \VerseOne{1}Căci acest Melchisedec, regele Salemului, preot al lui Dumnezeu Cel Preaînalt, care a întâlnit pe Avraam la întoarcerea de la măcelul regilor și l-a binecuvântat,

\VS{2}căruia Avraam i-a împărțit și a zecea parte din toate (fiind mai întâi, prin interpretare, “rege al dreptății”, apoi și “rege al Salemului”, adică “rege al păcii”,

\VS{3}fără tată, fără mamă, fără genealogie, fără început de zile și fără sfârșit de viață, ci făcut asemenea Fiului lui Dumnezeu), rămâne preot în permanență.


\par }{\PP \VS{4}Și gândește-te cât de mare era acest om, căruia chiar și patriarhul Avraam i-a dat o zecime din prada cea mai bună.

\VS{5}Într-adevăr, cei dintre fiii lui Levi care primesc slujba de preot au poruncă să ia zeciuiala de la popor, conform legii, adică de la frații lor, deși aceștia au ieșit din trupul lui Avraam,

\VS{6}dar cel a cărui genealogie nu se numără de la ei a acceptat zeciuiala de la Avraam și a binecuvântat pe cel care are făgăduințele.

\VS{7}Dar, fără nicio contestație, cel mai mic este binecuvântat de cel mai mare.

\VS{8}Aici primesc zeciuială oamenii care mor, dar acolo primește zeciuială unul despre care s-a mărturisit că trăiește.

\VS{9}Putem spune că, prin Avraam, chiar și Levi, care primește zeciuiala, a plătit zeciuiala,

\VS{10}căci el era încă în trupul tatălui său când Melchisedec l-a întâlnit.


\par }{\PP \VS{11}Și dacă desăvârșirea a fost prin preoția levitică, căci sub ea a primit poporul Legea, ce nevoie mai era ca un alt preot să se ridice după rânduiala lui Melchisedec, și să nu fie chemat după rânduiala lui Aaron?

\VS{12}Căci, fiind schimbată preoția, este necesar să se facă o schimbare și în lege.

\VS{13}Căci cel despre care se spun aceste lucruri aparține unei alte seminții, din care nimeni nu a oficiat la altar.

\VS{14}Căci este evident că Domnul nostru a ieșit din Iuda, trib despre care Moise nu a vorbit nimic cu privire la preoție.

\VS{15}Acest lucru este încă și mai evident dacă, după asemănarea lui Melchisedec, se ridică un alt preot,

\VS{16}care a fost făcut nu după legea unei porunci trupești, ci după puterea unei vieți fără sfârșit;

\VS{17}căci este mărturisit,

\par }{\Q “Ești preot pentru totdeauna,

\par }{\QB conform ordinului lui Melchisedec.”


\par }{\PP \VS{18}Căci este o anulare a unei porunci anterioare, din cauza slăbiciunii și a lipsei ei de folos

\VS{19}(căci Legea nu a făcut nimic desăvârșit), și o introducere a unei speranțe mai bune, prin care ne apropiem de Dumnezeu.

\VS{20}Întrucât nu a fost făcut preot fără jurământ

\VS{21}(căci, într-adevăr, au fost făcuți preoți fără jurământ), ci el cu jurământ, prin cel care spune despre el,

\par }{\Q “Domnul a jurat și nu se va răzgândi,

\par }{\QB 'Ești preot pentru totdeauna,

\par }{\QB conform ordinului lui Melchisedec.””


\par }{\PP \VS{22}Prin aceasta, Isus a devenit garanția unui legământ mai bun.


\par }{\PP \VS{23}Mulți au fost făcuți preoți, pentru că moartea i-a împiedicat să continue.

\VS{24}Dar el, pentru că trăiește veșnic, are o preoție neschimbată.

\VS{25}De aceea și el este în măsură să salveze până la capăt pe cei care se apropie de Dumnezeu prin el, întrucât el trăiește în veac ca să mijlocească pentru ei.


\par }{\PP \VS{26}Căci un astfel de mare preot se cuvenea pentru noi: sfânt, fără vină, fără pată, despărțit de păcătoși și mai înalt decât cerurile;

\VS{27}care nu are nevoie, ca acei mari preoți, să aducă zilnic jertfe, mai întâi pentru păcatele sale și apoi pentru păcatele poporului. Căci el a făcut acest lucru o dată pentru totdeauna, când s-a jertfit pe sine însuși.

\VS{28}Căci legea numește ca mari preoți pe oameni care au slăbiciuni, dar cuvântul jurământului, care a venit după lege, numește pentru totdeauna un Fiu care a fost desăvârșit.



\par }\Chap{8}{\PP \VerseOne{1}Or, în cele ce spunem, principalul este următorul lucru: avem un astfel de mare preot, care a șezut la dreapta tronului Maiestății în ceruri,

\VS{2}un slujitor al sanctuarului și al cortului adevărat, pe care l-a ridicat Domnul, nu omul.

\VS{3}Căci orice mare preot este desemnat să ofere atât daruri, cât și sacrificii. Prin urmare, este necesar ca și acest mare preot să aibă ceva de oferit.

\VS{4}Căci, dacă ar fi pe pământ, nu ar fi deloc preot, întrucât există preoți care oferă daruri conform Legii,

\VS{5}care slujesc ca o copie și o umbră a lucrurilor cerești, așa cum a fost avertizat Moise de către Dumnezeu când era pe cale să facă cortul, căci i-a spus: “Vezi, să faci totul după modelul care ți-a fost arătat pe munte”.

\VS{6}Dar acum a obținut o slujbă mai excelentă, prin faptul că este și mijlocitor al unui legământ mai bun, care, pe baza unor promisiuni mai bune, a fost dat ca lege.


\par }{\PP \VS{7}Căci dacă primul legământ ar fi fost fără cusur, nu s-ar fi căutat loc pentru al doilea.

\VS{8}Căci, găsind greșeli la ei, a spus,

\par }{\Q “Iată, zilele vin”, zice Domnul,

\par }{\QB “că voi face un legământ nou cu casa lui Israel și cu casa lui Iuda;


\par }{\Q \VS{9}nu după legământul pe care l-am făcut cu părinții lor.

\par }{\QB în ziua în care i-am luat de mână ca să-i scot din țara Egiptului;

\par }{\Q pentru că nu au continuat în legământul Meu,

\par }{\QB și nu le-am luat în seamă”, zice Domnul.


\par }{\Q \VS{10}“Căci acesta este legământul pe care-l voi face cu casa lui Israel.

\par }{\QB după acele zile”, zice Domnul:

\par }{\Q “Voi pune legile mele în mintea lor;

\par }{\QB Le voi scrie și pe inima lor.

\par }{\Q Eu voi fi Dumnezeul lor,

\par }{\QB și ei vor fi poporul Meu.


\par }{\Q \VS{11}Nu vor învăța pe fiecare dintre ei să-și învețe concetățeanul.

\par }{\QB și fiecare pe fratele său, zicând: “Cunoaște pe Domnul”.

\par }{\QB căci toți mă vor cunoaște,

\par }{\QB de la cel mai mic la cel mai mare.


\par }{\Q \VS{12}Căci voi fi milostiv cu nedreptatea lor.

\par }{\QB Nu-mi voi mai aduce aminte de păcatele lor și de faptele lor nelegiuite.”


\par }{\PP \VS{13}Când zice: “Un legământ nou”, îl face caduc pe cel dintâi. Dar ceea ce devine învechit și îmbătrânește este aproape de dispariție.



\par }\Chap{9}{\PP \VerseOne{1}Dar chiar și primul legământ avea rânduieli de serviciu divin și un sanctuar pământesc.

\VS{2}Căci era pregătit un cort. În prima parte se aflau sfeșnicul, masa și pâinile de arătat, care se numește Locul Sfânt.

\VS{3}După cel de-al doilea văl era Cortul care se numește Sfânta Sfintelor,

\VS{4}având un altar de aur pentru tămâie și chivotul legământului acoperit de aur pe toate părțile, în care se afla un vas de aur care ținea mana, toiagul lui Aaron care înmugurește și tablele legământului;

\VS{5}iar deasupra lui heruvimi de glorie care umbreau capacul milei, lucruri despre care nu putem vorbi acum în detaliu.


\par }{\PP \VS{6}Și, după ce aceste lucruri au fost astfel pregătite, preoții intră neîncetat în primul cort, îndeplinind slujbele,

\VS{7}dar în al doilea intră numai marele preot, o dată pe an, nu fără sânge, pe care îl aduce pentru el însuși și pentru greșelile poporului.

\VS{8}Duhul Sfânt indică acest lucru: că drumul spre Locul Sfânt nu fusese încă dezvăluit cât timp primul cortul era încă în picioare.

\VS{9}Acesta este un simbol al epocii actuale, în care se oferă daruri și sacrificii care, în ceea ce privește conștiința, nu sunt capabile să îl facă pe cel care se închină perfect,

\VS{10}fiind doar (cu mâncăruri și băuturi și diverse spălări) rânduieli trupești, impuse până la o vreme de reformă.


\par }{\PP \VS{11}Dar Hristos, venind ca Mare Preot al bunurilor viitoare, prin cortul cel mai mare și mai desăvârșit, care nu este făcut de mână, adică nu este din creația aceasta,

\VS{12}și nici prin sângele țapilor și vițeilor, ci prin propriul sânge, a intrat o dată pentru totdeauna în Locul Sfânt, obținând răscumpărarea veșnică.

\VS{13}Căci, dacă sângele țapilor și al taurilor și cenușa unei vițele care stropește pe cei care au fost spurcați sfințesc pentru curățirea cărnii,

\VS{14}cu cât mai mult sângele lui Cristos, care, prin Duhul etern, s-a oferit pe sine însuși fără cusur lui Dumnezeu, va curăța conștiința voastră de fapte moarte pentru a sluji Dumnezeului viu?

\VS{15}De aceea, el este mijlocitorul unui nou legământ, întrucât a avut loc o moarte pentru răscumpărarea fărădelegilor care au fost sub primul legământ, pentru ca cei care au fost chemați să primească promisiunea moștenirii veșnice.

\VS{16}Căci acolo unde este un testament, trebuie neapărat să fie și moartea celui care l-a făcut.

\VS{17}Căci testamentul este în vigoare acolo unde a avut loc moartea, căci el nu este niciodată în vigoare cât timp trăiește cel care l-a făcut.

\VS{18}De aceea, nici măcar primul legământ nu a fost consacrat fără sânge.

\VS{19}Căci, după ce Moise a rostit fiecare poruncă, conform Legii, către tot poporul, a luat sângele vițeilor și al țapilor, cu apă, lână stacojie și isop, și a stropit atât cartea însăși, cât și tot poporul,

\VS{20}spunând: “Acesta este sângele legământului pe care vi l-a poruncit Dumnezeu”.


\par }{\PP \VS{21}Și a stropit cu sânge și cortul și toate uneltele slujbei.

\VS{22}Potrivit legii, aproape totul se curăță cu sânge și, în afară de vărsarea de sânge, nu există iertare.


\par }{\PP \VS{23}De aceea era necesar ca prin acestea să fie curățate copiile lucrurilor din ceruri, iar lucrurile cerești să fie curățate cu jertfe mai bune decât acestea.

\VS{24}Căci Hristos nu a intrat în locurile sfinte făcute de mâini, care sunt reprezentări ale celor adevărate, ci în cerul însuși, ca să se înfățișeze acum în fața lui Dumnezeu pentru noi;

\VS{25}și nici că ar trebui să se aducă pe Sine însuși adeseori, așa cum intră marele preot în locul sfânt an de an, cu sânge care nu este al Lui,

\VS{26}altfel trebuie să fi suferit adeseori de la întemeierea lumii. Dar acum, o singură dată, la sfârșitul veacurilor, el s-a arătat pentru a înlătura păcatul prin jertfa lui însuși.

\VS{27}În măsura în care oamenilor le este rânduit să moară o singură dată, iar după aceasta, judecata,

\VS{28}tot așa și Hristos, după ce a fost oferit o singură dată pentru a purta păcatele multora, se va arăta a doua oară, nu pentru a se ocupa de păcat, ci pentru a-i salva pe cei care îl așteaptă cu nerăbdare.



\par }\Chap{10}{\PP \VerseOne{1}Fiindcă Legea, care este o umbră a bunurilor viitoare, și nu chipul însuși al lucrurilor, nu poate niciodată, cu aceleași jertfe pe care le aduce an de an, pe care le aduce neîncetat, să desăvârșească pe cei ce se apropie.

\VS{2}Altfel, nu ar fi încetat ele să mai fie oferite, pentru că închinătorii, după ce ar fi fost curățați o dată, nu ar mai fi avut conștiința păcatelor?

\VS{3}Dar în aceste sacrificii există o reamintire anuală a păcatelor.

\VS{4}Căci este imposibil ca sângele taurilor și al țapilor să înlăture păcatele.

\VS{5}De aceea, atunci când vine în lume, spune: “Nu, nu, nu!

\par }{\Q “Tu nu ai dorit sacrificii și jertfe,

\par }{\QB dar tu mi-ai pregătit un trup.


\par }{\Q \VS{6}Nu-ți plăceau arderile de tot și jertfele pentru păcat.


\par }{\QB \VS{7}Atunci am zis: “Iată, am venit (în sulul cărții este scris despre mine).

\par }{\QB ca să fac voia Ta, Dumnezeule.””


\par }{\PP \VS{8}Mai înainte, zicând: “Jertfe și daruri și arderi de tot și jertfe pentru păcat nu ai dorit, nici nu ți-a plăcut” (cele ce se aduc după lege),

\VS{9}apoi a zis: “Iată, am venit să fac voia Ta.” El le înlătură pe cele dintâi, ca să le stabilească pe cele de a doua,

\VS{10}prin a căror voință am fost sfințiți prin jertfa trupului lui Isus Cristos o dată pentru totdeauna.


\par }{\PP \VS{11}Orice preot stă zi de zi, slujind și aducând mereu aceleași jertfe, care nu pot niciodată să ridice păcatele,

\VS{12}dar El, după ce a adus o singură jertfă pentru păcate pentru totdeauna, a șezut la dreapta lui Dumnezeu,

\VS{13}și de atunci așteaptă până când vrăjmașii Lui vor fi așezate sub picioarele Lui.

\VS{14}Căci printr-o singură jertfă a desăvârșit pentru totdeauna pe cei care sunt sfințiți.

\VS{15}Și Duhul Sfânt ne mărturisește, de asemenea, că, după ce a spus


\par }{\Q \VS{16}“Acesta este legământul pe care-l voi face cu ei.

\par }{\QB după acele zile”, zice Domnul,

\par }{\Q “Voi pune legile Mele în inima lor,

\par }{\QB Le voi scrie și în mintea lor;”

\par }{\PP apoi spune,


\par }{\Q \VS{17}“Nu-mi voi mai aduce aminte de păcatele lor și de nelegiuirile lor.”


\par }{\PP \VS{18}Și unde este iertarea acestora, nu mai este nici o jertfă pentru păcat.


\par }{\PP \VS{19}Așadar, fraților, având îndrăzneala de a intra în locul sfânt prin sângele lui Isus,

\VS{20}pe calea pe care El a deschis-o pentru noi, o cale nouă și vie, prin văl, adică prin trupul Său,

\VS{21}și având un mare preot peste casa lui Dumnezeu,

\VS{22}să ne apropiem cu inimă sinceră, în plinătatea credinței, având inimile stropite de conștiința rea și trupul spălat cu apă curată,

\VS{23}să ținem neclintit mărturisirea speranței noastre, căci Cel ce a făgăduit este credincios.


\par }{\PP \VS{24}Să ne gândim cum să ne îndemnăm unii pe alții la iubire și la fapte bune,

\VS{25}fără a renunța la adunarea noastră, cum obișnuiesc unii, ci îndemnându-ne unii pe alții, și cu atât mai mult cu cât vedeți că se apropie Ziua.


\par }{\PP \VS{26}Căci, dacă păcătuim cu bună știință după ce am primit cunoștința adevărului, nu mai rămâne nici o jertfă pentru păcate,

\VS{27}ci o așteptare înfricoșătoare a judecății și o înflăcărare de foc care va mistui pe cei potrivnici.

\VS{28}Un om care nu respectă legea lui Moise moare fără milă, pe cuvântul a doi sau trei martori.

\VS{29}Cu cât mai rea pedeapsă credeți că va fi judecat vrednic cel care a călcat în picioare pe Fiul lui Dumnezeu, care a socotit ca fiind un lucru profanat sângele legământului cu care a fost sfințit și care a insultat Duhul harului?

\VS{30}Căci noi îl cunoaștem pe cel care a spus: “Răzbunarea îmi aparține. Eu voi răsplăti”, zice Domnul. Din nou: “Domnul va judeca poporul său”.

\VS{31}Este un lucru înfricoșător să cazi în mâinile Dumnezeului cel viu.


\par }{\PP \VS{32}Dar adu-ți aminte de zilele de mai înainte, când, după ce ai fost luminat, ai dus o mare luptă cu suferințe:

\VS{33}pe de o parte, fiind expuși atât la ocară, cât și la asupririri, iar pe de altă parte, devenind părtași cu cei care erau tratați astfel.

\VS{34}Căci amândoi ați avut milă de mine în lanțurile mele și ați acceptat cu bucurie jefuirea averilor voastre, știind că aveți pentru voi înșivă o avere mai bună și durabilă în ceruri.

\VS{35}De aceea, nu aruncați la gunoi îndrăzneala voastră, care are o mare răsplată.

\VS{36}Căci aveți nevoie de rezistență pentru ca, după ce ați făcut voia lui Dumnezeu, să primiți făgăduința.


\par }{\Q \VS{37}“Peste foarte puțin timp,

\par }{\QB cel ce vine va veni și nu va aștepta.


\par }{\Q \VS{38}Dar cel neprihănit va trăi prin credință.

\par }{\QB Dacă se retrage, sufletul meu nu are nici o plăcere în el.”


\par }{\PP \VS{39}Dar noi nu suntem dintre cei ce se retrag spre pieire, ci dintre cei ce au credință pentru mântuirea sufletului.



\par }\Chap{11}{\PP \VerseOne{1}Or, credința este o asigurare a lucrurilor nădăjduite, o dovadă a lucrurilor care nu se văd.

\VS{2}Căci, prin aceasta, bătrânii au obținut aprobarea.

\VS{3}Prin credință înțelegem că universul a fost alcătuit prin Cuvântul lui Dumnezeu, astfel încât ceea ce se vede nu a fost făcut din lucruri vizibile.


\par }{\PP \VS{4}Prin credință, Abel a adus lui Dumnezeu o jertfă mai bună decât Cain, prin care i s-a dat mărturie că era neprihănit, Dumnezeu mărturisind despre darurile sale; și prin ea, el, mort fiind, vorbește încă.


\par }{\PP \VS{5}Prin credință, Enoh a fost luat, ca să nu vadă moartea, și nu a fost găsit, pentru că Dumnezeu l-a mutat. Căci i s-a dat mărturie că înainte de translarea lui a fost bineplăcut lui Dumnezeu.

\VS{6}Fără credință este imposibil să-i fii bineplăcut, căci cine se apropie de Dumnezeu trebuie să creadă că el există și că este răsplătitorul celor care îl caută.


\par }{\PP \VS{7}Prin credință, Noe, fiind avertizat despre lucruri care nu se vedeau încă, a pregătit, cu frică evlavioasă, o corabie pentru salvarea casei sale, prin care a condamnat lumea și a devenit moștenitor al dreptății care este după credință.


\par }{\PP \VS{8}Prin credință, Avraam, când a fost chemat, a ascultat și a ieșit la locul pe care avea să-l primească de moștenire. A ieșit, fără să știe unde se duce.

\VS{9}Prin credință a trăit ca un străin în țara făgăduinței, ca într-o țară care nu era a lui, locuind în corturi cu Isaac și Iacov, moștenitorii împreună cu el ai aceleiași făgăduințe.

\VS{10}Căci el aștepta cetatea care are temelii, al cărei ziditor și făcător este Dumnezeu.


\par }{\PP \VS{11}Prin credință, chiar și Sara a primit puterea de a rămâne însărcinată și a născut un copil la o vârstă înaintată, pentru că a socotit credincios pe Cel ce făgăduise.

\VS{12}De aceea, câte stele sunt în mulțime ca stelele cerului și nenumărate ca nisipul care este pe malul mării, au fost generate de un singur om, și el ca și mort.


\par }{\PP \VS{13}Toți aceștia au murit cu credință, nu după ce au primit făgăduințele, ci după ce le-au văzut și le-au îmbrățișat de departe, și după ce au mărturisit că sunt străini și pelerini pe pământ.

\VS{14}Căci cei care spun astfel de lucruri arată clar că își caută o țară a lor.

\VS{15}Dacă într-adevăr s-ar fi gândit la țara din care au plecat, ar fi avut destul timp să se întoarcă.

\VS{16}Dar acum ei doresc o țară mai bună, adică una cerească. De aceea Dumnezeu nu se rușinează de ei, ca să fie numit Dumnezeul lor, căci le-a pregătit o cetate.


\par }{\PP \VS{17}Avraam, fiind pus la încercare, a jertfit pe Isaac prin credință. Da, cel care primise cu bucurie făgăduințele își oferea singurul său fiu născut,

\VS{18}căruia i s-a spus: “Descendența ta va fi socotită ca din Isaac”,

\VS{19}concluzionând că Dumnezeu este în stare să învie chiar și din morți. La figurat vorbind, l-a și primit înapoi din morți.


\par }{\PP \VS{20}Prin credință, Isaac a binecuvântat pe Iacov și pe Esau cu privire la cele viitoare.


\par }{\PP \VS{21}Prin credință, Iacov, pe când era pe moarte, a binecuvântat pe fiecare dintre fiii lui Iosif și s-a închinat, sprijinindu-se în vârful toiagului său.


\par }{\PP \VS{22}Prin credință, Iosif, când i s-a apropiat sfârșitul, a anunțat plecarea copiilor lui Israel și a dat instrucțiuni cu privire la oasele lui.


\par }{\PP \VS{23}Prin credință, Moise, când s-a născut, a fost ascuns trei luni de părinții săi, pentru că au văzut că era un copil frumos, și nu s-au temut de porunca împăratului.


\par }{\PP \VS{24}Prin credință, Moise, după ce a crescut, a refuzat să fie numit fiul fiicei lui Faraon,

\VS{25}alegând mai degrabă să se împartă cu poporul lui Dumnezeu decât să se bucure pentru o vreme de plăcerile păcatului,

\VS{26}considerând că ocara lui Hristos este o bogăție mai mare decât comorile Egiptului, pentru că se aștepta la răsplată.

\VS{27}Prin credință a părăsit Egiptul, fără să se teamă de mânia regelui, căci a îndurat, ca și cum ar fi văzut pe Cel care este invizibil.

\VS{28}Prin credință a păzit Paștele și stropirea sângelui, pentru ca nimicitorul primilor născuți să nu se atingă de ei.


\par }{\PP \VS{29}Prin credință au trecut prin Marea Roșie ca pe uscat. Când egiptenii au încercat să facă acest lucru, au fost înghițiți.


\par }{\PP \VS{30}Prin credință au căzut zidurile Ierihonului, după ce au fost înconjurate timp de șapte zile.


\par }{\PP \VS{31}Prin credință, prostituata Rahab nu a pierit împreună cu cei neascultători, ci a primit pe spioni în pace.


\par }{\PP \VS{32}Ce să mai spun? Căci nu aș mai avea timp dacă aș povesti despre Ghedeon, Barac, Samson, Iefta, David, Samuel și proorocii,

\VS{33}care, prin credință, au supus împărății, au făcut dreptate, au obținut făgăduințe, au astupat gurile leilor,

\VS{34}au stins puterea focului, au scăpat de ascuțișul sabiei, din slăbiciune s-au întărit, au devenit puternici în război și au făcut să fugă oștirile străine.

\VS{35}Femeile și-au primit morții prin înviere. Alții au fost torturați, neacceptând eliberarea, pentru a obține o înviere mai bună.

\VS{36}Alții au fost încercați prin batjocură și biciuire, da, mai mult, prin legături și închisoare.

\VS{37}Au fost uciși cu pietre. Au fost tăiați cu ferăstrăul. Au fost ispitiți. Au fost uciși cu sabia. Au umblat în piei de oaie și în piei de capră; fiind lipsiți, chinuiți, maltratați —

\VS{38}de care lumea nu era demnă — rătăcind prin deșerturi, munți, peșteri și grote și prin găurile pământului.


\par }{\PP \VS{39}Toți aceștia, după ce au fost lăudați pentru credința lor, n-au primit făgăduința,

\VS{40}Dumnezeu a pregătit ceva mai bun în privința noastră, astfel încât, în afară de noi, ei n-ar fi fost desăvârșiți.



\par }\Chap{12}{\PP \VerseOne{1}De aceea și noi, fiindcă suntem înconjurați de un nor atât de mare de martori, să ne lepădăm de orice greutate și de păcatul care ne încolțește atât de ușor, și să alergăm cu stăruință la cursa care ne este pusă înainte,

\VS{2}privind la Isus, autorul și desăvârșitorul credinței, care, pentru bucuria care i-a fost pusă înainte, a îndurat crucea, disprețuind rușinea ei, și a șezut la dreapta tronului lui Dumnezeu.


\par }{\PP \VS{3}Căci gândiți-vă la El, care a îndurat atâtea împotriviri ale păcătoșilor împotriva lui însuși, ca să nu vă obosiți, să nu vă pierdeți curajul sufletelor voastre.

\VS{4}Voi nu v-ați împotrivit încă până la sânge, luptând împotriva păcatului.

\VS{5}Ați uitat îndemnul care raționează cu voi ca și cu copiii,

\par }{\Q “Fiule, nu lua în seamă pedeapsa Domnului,

\par }{\QB nici să nu leșinați când sunteți mustrați de el;


\par }{\QB \VS{6}Căci pe cine iubește Domnul, El îl pedepsește,

\par }{\QB și pedepsește orice fiu pe care îl primește.”


\par }{\PP \VS{7}Pentru că pentru disciplinarea voastră sunteți înduplecați. Dumnezeu se poartă cu voi ca și cu copiii, căci ce fiu este acela pe care tatăl său nu-l pedepsește?

\VS{8}Dar dacă sunteți lipsiți de disciplină, de care toți au fost făcuți părtași, atunci sunteți nelegitimi și nu copii.

\VS{9}Mai mult, noi am avut părinți din carnea noastră care să ne pedepsească, și le-am dat respect. Nu ar trebui mult mai degrabă să ne supunem Tatălui spiritelor și să trăim?

\VS{10}Căci ei, într-adevăr, pentru câteva zile ne-au disciplinat cum li s-a părut lor bine, dar el pentru folosul nostru, ca să fim părtași la sfințenia lui.

\VS{11}Orice pedeapsă pare pentru moment nu veselă, ci dureroasă; dar, după aceea, ea dă roadele pașnice ale neprihănirii celor care au fost instruiți prin ea.

\VS{12}De aceea, ridicați mâinile care atârnă și genunchii slabi,

\VS{13}și faceți cărări drepte pentru picioarele voastre, ca ceea ce este șchiop să nu fie dislocat, ci mai degrabă să fie vindecat.


\par }{\PP \VS{14}Urmăriți pacea cu toți oamenii și sfințirea, fără de care nimeni nu va vedea pe Domnul,

\VS{15}privind cu atenție ca nu cumva să fie cineva care să fie lipsit de harul lui Dumnezeu, ca nu cumva să vă tulbure vreo rădăcină de amărăciune care să răsară și mulți să fie spurcați de ea,

\VS{16}ca nu cumva să fie vreun desfrânat sau vreun profanator, ca Esau, care și-a vândut dreptul de întâi-născut pentru o masă.

\VS{17}Căci știți că și atunci când a dorit ulterior să moștenească binecuvântarea, a fost respins, pentru că nu a găsit loc de răzgândire, deși a căutat-o cu sârguință și cu lacrimi.


\par }{\PP \VS{18}Căci nu ați venit la un munte care putea fi atins și care ardea cu foc, la negură, întuneric, furtună,

\VS{19}la sunetul de trâmbiță și la glasul cuvintelor, pe care cei care le-au auzit rugau să nu li se mai spună nici un cuvânt,

\VS{20}căci nu puteau să suporte ceea ce era poruncit: “Dacă un animal se atinge de munte, să fie ucis cu pietre”.

\VS{21}Atât de înfricoșătoare era înfățișarea, încât Moise a spus: “Sunt îngrozit și tremur”.


\par }{\PP \VS{22}Dar voi ați venit la muntele Sionului și la cetatea Dumnezeului celui viu, Ierusalimul cel ceresc, la mulțimile nenumărate de îngeri,

\VS{23}la adunarea de sărbătoare și la adunarea întâilor născuți care sunt înscriși în ceruri, la Dumnezeu, Judecătorul tuturor, la duhurile drepților desăvârșiți,

\VS{24}la Isus, mijlocitorul unui legământ nou, și la sângele stropirii, care este mai bun decât cel al lui Abel.


\par }{\PP \VS{25}Vezi să nu refuzi pe cel ce vorbește. Căci dacă ei nu au scăpat când l-au refuzat pe cel care avertiza pe pământ, cu atât mai mult nu vom scăpa noi, care ne întoarcem de la cel care avertizează din ceruri,

\VS{26}al cărui glas a zguduit atunci pământul, dar acum a promis: “Încă o dată voi mai zgudui nu numai pământul, ci și cerurile”.

\VS{27}Această expresie: “Încă o dată” semnifică îndepărtarea lucrurilor care sunt zguduite, ca și a lucrurilor care au fost făcute, pentru ca cele care nu sunt zguduite să rămână.

\VS{28}Prin urmare, primind o Împărăție care nu poate fi zdruncinată, să avem har, prin care să Îi slujim lui Dumnezeu în mod acceptabil, cu reverență și respect,

\VS{29}căci Dumnezeul nostru este un foc mistuitor.



\par }\Chap{13}{\PP \VerseOne{1}Să continue dragostea frățească.

\VS{2}Nu uitați să fiți ospitalieri cu străinii, pentru că, făcând astfel, unii au găzduit îngeri fără să știe.

\VS{3}Aduceți-vă aminte de cei care sunt în legături, ca și cum ați fi legați cu ei, și de cei maltratați, întrucât și voi sunteți în trup.

\VS{4}Căsătoria să fie ținută în cinste între toți și patul să fie neîntinat; dar Dumnezeu îi va judeca pe cei imorali și pe adulteri.


\par }{\PP \VS{5}Nu vă lăsați de iubirea de bani, mulțumiți-vă cu ce aveți, căci El a zis: “Nu vă voi lăsa și nu vă voi părăsi nicidecum”.

\VS{6}astfel încât, cu mult curaj, să spunem

\par }{\Q “Domnul este ajutorul meu. Nu mă voi teme.

\par }{\QB Ce-mi poate face omul?”


\par }{\PP \VS{7}Aduceți-vă aminte de conducătorii voștri, care v-au vestit Cuvântul lui Dumnezeu, și, ținând seama de rezultatele faptelor lor, imitați-le credința.

\VS{8}Isus Hristos este același ieri, astăzi și în veci.

\VS{9}Nu vă lăsați purtați de învățături diferite și ciudate, căci este bine ca inima să fie întărită prin har, nu prin mâncăruri, de care nu au beneficiat cei care s-au ocupat astfel.


\par }{\PP \VS{10}Avem un altar din care nu au voie să mănânce cei ce slujesc la sfântul cort.

\VS{11}Căci trupurile acelor animale, al căror sânge este adus în sfântul locaș de către marele preot ca jertfă pentru păcat, sunt arse în afara taberei.

\VS{12}De aceea și Isus, ca să sfințească poporul prin sângele Său, a suferit în afara porții.

\VS{13}Să mergem deci la el în afara taberei, purtându-i ocara.

\VS{14}Căci noi nu avem aici o cetate durabilă, ci căutăm pe cea care va veni.

\VS{15}Prin el, așadar, să aducem permanent prin el o jertfă de laudă lui Dumnezeu, adică rodul buzelor care proclamă credință numelui său.

\VS{16}Dar nu uitați să faceți binele și să împărțiți, căci cu astfel de sacrificii Dumnezeu este binevoitor.


\par }{\PP \VS{17}Ascultați de căpeteniile voastre și supuneți-vă lor, căci ei veghează pentru sufletele voastre, ca cei ce vor da socoteală, ca să facă acest lucru cu bucurie, nu cu gemete, căci aceasta ar fi nefolositor pentru voi.


\par }{\PP \VS{18}Rugați-vă pentru noi, pentru că suntem încredințați că avem o conștiință bună, dorind să trăim cu cinste în toate lucrurile.

\VS{19}Vă îndemn cu tărie să faceți acest lucru, pentru ca eu să vă fiu redat mai repede.


\par }{\PP \VS{20}Și Dumnezeul păcii, care a înviat din morți pe marele păstor al oilor cu sângele unui legământ veșnic, pe Domnul nostru Isus,

\VS{21}să vă facă desăvârșiți în orice lucrare bună, ca să faceți voia Lui, lucrând în voi ceea ce este plăcut înaintea Lui, prin Isus Hristos, căruia I se cuvine slava în vecii vecilor. Amin.


\par }{\PP \VS{22}Dar eu vă îndemn, fraților, să îndurați cuvântul de îndemn, căci v-am scris în puține cuvinte.

\VS{23}Să știți că a fost eliberat fratele nostru Timotei, cu care, dacă va veni curând, vă voi vedea.


\par }{\PP \VS{24}Salutați pe toți conducătorii voștri și pe toți sfinții. Italienii vă salută.


\par }{\PP \VS{25}Harul să fie cu voi toți. Amin.


\par }